%%=====================================================================
%%  3 微信小程序开发需求设计
%%=====================================================================
\section{微信小程序开发需求设计}

\subsection{用户需求分析}

PMP 刷题系统的用户对象是所有使用微信且有 PMP 刷题需求的用户，用户需求主要有以下几个方面：

\begin{enumerate}
  \item 用户可以选择每次答题的数量和章节；
  \item 用户可以进行在线快速答题；
  \item 用户答题完毕后可以查看答案解析；
  \item 用户可以查看历史记录，查看已做过题目的答案解析；
  \item 用户可将小程序发送至微信好友或微信群；
  \item 用户可以查看不同群的好友答题得分排名；
  \item 用户遇到问题可以进行留言反馈；
  \item 用户可以查看距离下一次考试的天数倒计时；
  \item 用户每日可进行一次运势占卜，占卜的结果可分享或保存在本地。
\end{enumerate}

\subsection{功能需求分析}

以用户需求为出发点，结合微信小程序开发的基本要求，本 PMP 刷题系统主要实现在线快速刷题、查看历史记录、查看群排名、幸运占卜、咨询反馈、分享等在线功能。具体要求如下。

快速刷题：用户能够选择每次答题数目和答题章节，并且答题完毕后才显示"提交答题"按钮。提交答卷后显示本次答题正确率，用户可以选择继续刷题或者查看本次答题解析。

查看历史记录：可以查看已答过的试卷信息，以红绿两色区分正确与否。点击对应的答题记录可以查看对应的题目解析。

查看群排行榜：可以查看已经分享过的微信群内的群成员答题总分以及群排名，并且不同的微信群中的排名可能不同。

咨询反馈：用户可以将在小程序使用过程中的问题，通过留言板反馈给后台。反馈的信息中除用户描述的问题，还会包含用户的设备信息。

幸运占卜：用户每日做满 20 题，可以进行一次幸运抽牌。包含洗牌、蜡烛移动、眨眼、音效播放等动画效果。

分享：用户可以将当日抽签的卡片、答题群排行榜得分分享至微信群，其他微信用户可以从分享的卡片进入小程序。同时，用户可以将当日抽取的卡片保存至手机相册，并可主动分享图片至微信朋友圈。
